% !PRINTED 451
\weberpage{451}
Now the point $\pi_0$ lies on the three principal axes of the substitutions of degree $2$, $\omega, \omega\chi, \omega\chi^2$. If there is yet a further substitution of degree $2$ by which $\pi_0$ is left unchanged, then $\pi_0$ must lie on the principal axis of this substitution.

But if we set $x_1 = x_2 = x_3$ in the equations of the principal axes [\S.\,118, (7)], each of these three equations is satisfied for only one value of $r$; hence $\pi_0$ cannot lie on more than three of the principal axes, and $\pi_0$ is six-fold.

From this the theorem:

7. \gesperrt{There are 42 six-fold poles, which can all be derived from one of them by substitutions of the group. Through each of these poles pass three principal axes of the group.}

The pole $\pi_0$ lies on none of the lines $T_{s,r}$. This is easily shown by setting $x_1 = x_2 = x_3$ in the expressions (4), and then putting in for $\alpha^2, \beta^2, \gamma^2$ their expressions in terms of $\varepsilon$. One then sees easily (with an application of the irreducibility of the cyclotomic equation) that none of these expressions vanishes.

\S.\,120.

\gesperrt{The Configuration of the Group $G_{168}$.}

The totality of the straight lines and points that we have considered in the preceding sections we shall call the \gesperrt{configuration of the group} $G_{168}$. The word here has the same meaning in which it has been used in recent times in geometry\footnote{The expression was introduced by \gesperrt{Reye}: \glqq{}Geometrie der Lage.\grqq{} Bd. I, 2. Aufl. (1876).}.

In what follows we describe this configuration somewhat more closely, without adding anything new, only in order to let the theorems of the last sections stand out more vividly and more clearly.

We have in this configuration:
\[
\begin{aligned}
&\text{21 lines } A \text{ (the principal axes)}, \\
&\text{21 points } a \text{ (the eight-fold poles)}, \\
&\text{28 lines } T, \\
&\text{28 points } t \text{ (the six-fold poles)}.
\end{aligned}
\]
